\begin{frame}
Dynamic Systems
 \begin{itemize}
 \item The time-dependent behavior of physical systems can often
 be modeled by differential equations.
 \item Such systems are often called \textit{\textbf{dynamic systems}} or
 \textit{\textbf{dynamical systems}}.
\end{itemize}
\end{frame}
\begin{frame}
Dynamic Systems
 \begin{itemize}
 \item Considering Newton's law $F=ma$ once again, we see that the
 force $F$ is the resultant force acting on the mass $m$.
 \item However, the resultant force is simply the vector sum of all the forces, both those
 internal to the system as well as the input forces that are external.
 \item The restoring force, $-ky$ in the simple oscillator
 problem, is a force internal to the system.  So we don't consider
 it an input \textit{per se}.
\end{itemize}
\end{frame}
\begin{frame}
Dynamic Systems
 \begin{itemize}
 \item The response or output of the system is motion, more specifically,
 accelerated motion given by $a$.
 \item Thus, the simple harmonic oscillator is a dynamic system, one in which there is only an internal force determining
 system behavior: $F=-ky$.
 \item An external forcing function may also be present, e.g.:
 \[\frac{d^2y}{dt^2}+\frac{k}{m}y=\sin\omega t\]
\end{itemize}
\end{frame}
\begin{frame}
Dynamic Systems
 \begin{itemize}
 \item The most general second-order equation for a mechanical
 system of this sort according to Newton's law is:
\begin{equation}\label{second ord}
  \frac{d^2y}{dt^2} =\frac{1}{m}F(y,\dot{y},t)
\end{equation}
 \item If we can separate the forcing function $F(y,\dot{y},t)$
 into a sum of two functions $F_1(y,\dot{y})$ and $F_2(t)$, then
 we can rearrange this equation into the input-output form of
 equation \eqref{diffeq}:
  \[\frac{d^2y}{dt^2}-\frac{1}{m}F_1(y,\dot{y})=\frac{1}{m}F_2(t)\]
 \end{itemize}
\end{frame}
\begin{frame}
Dynamic Systems
 \begin{itemize}
 \item In engineering applications, we are often interested not
 only in modeling the dynamics of the system, but also in
 measuring the output or behavior of the system.
 \item But the output is not always directly observable.
 \item Thus, in specifying dynamic systems, we often add a second
 equation, usually algebraic or functional in nature, that relates
 the output $y$ (and possibly the input $u$) to some measured or observed quantities.
  \[z = h(y)+g(u)\]
\end{itemize}
\end{frame}
\begin{frame}
Dynamic Systems: The Concept of State
 \begin{itemize}
 \item In keeping with the linear spirit of things, we will
 restrict ourselves for now to linear functions $h$.
 \item Also, we will now introduce the state space concept and
 define $z$ as a function of the state.
 \item Intuitively, the concept of state involves in some sense the
 instantaneous internal arrangement of a system.  In practical
 situations, the state variables may represent attributes of the
 system that we wish to know something about.
\end{itemize}
\end{frame}
\begin{frame}
Dynamic Systems: The Concept of State
 \begin{itemize}
 \item Mechanical systems like a moving particle have three simultaneous attributes:  position, velocity, and
 acceleration.
 \item Newton tells us that acceleration is often a function
 of position, velocity and perhaps time.
 \item However, as we just saw, a force that acts on a system as a function of
 time only should be considered an external forcing function,
 i.e., an input to the system.
\item Thus we are led to think of position and velocity as the
primary state variables of interest.
\end{itemize}
\end{frame}
\begin{frame}
State Space Representation
 \begin{itemize}
 \item A general $n$-th order linear differential equation can be
 expressed as a system of first-order differential equations.
 \item Ignoring the input $u$ for the moment, consider:
 \[p(D)y = D^ny+a_{n-1}D^{n-1}y+\cdots +a_1Dy+a_0y,\]
 where we have put the polynomial in normal form, i.e., such that
 its leading coefficient is 1.
\end{itemize}
\end{frame}
\begin{frame}
State Space Representation
 \begin{itemize}
 \item If we consider $y$ and its first $n-1$ derivatives as the
 states of this system, then we have $n$ first order equations.
 Setting $x_1=y$, we write
\begin{align*}
  x_1(t) & = y(t) \\
  x_2(t) & = Dy(t) \\
    \cdots \\
  x_n(t) & = D^{n-1}y(t).
\end{align*}
\item The functions $x_i(t)$ are called the \textit{\textbf{state variables}}.
\end{itemize}
\end{frame}
\begin{frame}
State Space Representation
 \begin{itemize}
 \item Taking the derivative of these $n$ state variables and making use of the
 original differential equation yields:
\begin{align*}
  Dx_1(t) & = x_2(t) \\
  Dx_2(t) & = x_3(t) \\
    \cdots \\
  Dx_n(t) & =-a_{n-1}x_n(t)-\cdots - a_0x_1(t).
\end{align*}
 \item  These $n$ equations can be written more compactly in vector matrix
 form:  $Dx(t)=Ax(t)$ or $\dot{x}(t)=Ax(t)$.
\end{itemize}
\end{frame}
\begin{frame}
State Space Representation
 \begin{itemize}
 \item In the equation $\dot{x}(t)=Ax(t)$, we have:
 \[x(t)=\begin{bmatrix}x_1(t)\\x_2(t)\\\vdots \\
 x_n(t)\end{bmatrix}\]
 and
 \[A=\begin{bmatrix}
 0 & 1 & 0 &\cdots & 0\\
 0 & 0 & 1 & \cdots & 0\\
 \vdots & & \ddots & & \vdots \\
 0 & 0  & 0 & 0 & 1\\
 -a_{0} &-a_{1} & \cdots & -a_{n-2} & -a_{n-1}
 \end{bmatrix},\]
 the \textit{\textbf{companion matrix}} of the differential
 equation.
\end{itemize}
\end{frame}
\begin{frame}
State Space Representation
 \begin{itemize}
 \item If there is an input function $u$ and it is present via its
 derivatives as in equation \eqref{diffeq}, then we can write:
\begin{equation}\label{statedyn}
  \dot{x}=Ax+Bu
\end{equation}
forming $B$ in much the same way we formed $A$, and suppressing
the time variable $t$.
 \item Equation \eqref{statedyn} expresses the system's \textit{\textbf{state dynamics}}.
 \item Later we will include white noise as an input.  However,
 that will make the differential equations stochastic, and we are not
 yet ready for that.
\end{itemize}
\end{frame}
\begin{frame}
State Space Representation
 \begin{itemize}
 \item The output model to this we obtain one
 additional vector-matrix equation:
\begin{equation}\label{measdyn1}
  y = Cx + Du
\end{equation}
where the matrix $D$ represents a direct feedthrough of the input to the output.  
\item The output equation is derived directly from the $b_i$ coefficients in the original matrix differential equation.
\end{itemize}
\end{frame}
\begin{frame}
State Space Representation
 \begin{itemize}
  \item Kalman filter theory generally assumes $D=0$ in the output equation and uses a measurement model instead:
\begin{equation}\label{measdyn2}
  z = Hx
\end{equation}
 \item Note that $H$ implicitly includes $C$ as a factor.
 \item Note that we have not placed any restrictions on the size of
 the measurement column vector $z$.  As long as $H$ has $n$ columns it can
 have as many rows as desired.
\end{itemize}
\end{frame}
\begin{frame}
State Space Representation: Example
 \begin{itemize}
 \item Let's consider once again the simple harmonic oscillator:
 \[m\frac{d^2y}{dt^2}+ky=0\]
 and convert this to state space form:
\begin{align*}
  x_1 & =y \\
  x_2 & =Dy
\end{align*}
and $\dot{x}_2 = - \frac{k}{m}x_1+ 0\cdot x_2 $.
 \item Therefore, in matrix and vector form:
 \[
 \begin{bmatrix}
 \dot{x}_1\\
 \dot{x}_2
 \end{bmatrix}
 =
\begin{bmatrix}
 0 & 1 \\
 -k/m & 0
 \end{bmatrix}
 \begin{bmatrix}
 x_1\\
 x_2
 \end{bmatrix}.
 \]
\end{itemize}
\end{frame}
\begin{frame}
\begin{itemize}
 \item For a measurement model, we could consider a camera with a strobe flash that periodically records the instantaneous position of the sliding mass.
 \item In this case, only the position would be observable and we would write the measurement or observation model as:
 \[
 z =
\begin{bmatrix}
 1 & 0
 \end{bmatrix}
 \begin{bmatrix}
 x_1\\
 x_2
 \end{bmatrix}.
 \]
(Remember that we are ignoring both process and measurement noise.)
\end{itemize}
\end{frame}

\begin{frame}
State Space Representation: Matlab / Simulink Example
\begin{itemize}
  \item Four parameters $A,B,C,D$ from Matlab workspace
  \item Measurement model added external to Simulink's state space block
\end{itemize} 
\end{frame}

\begin{frame}
State Space Representation:  Summary
 \begin{itemize}
 \item The state variables of the system are usually provided by
 the 0th through $n-1$st derivatives.
 \item Equivalently, they are the outputs of the $n$ integrators of
 the system.
 \item They represent the internal coupling of the system.
 \item Any similarity transformation of the state vector yields an
 equivalent state space representation.  This is nothing more than
 changing coordinate systems (and is sometimes very useful).
\end{itemize}
\end{frame}
